\chapter{多元函数微分学}\label{ch:6}
\begin{tcolorbox}[title=章首知识地图,colframe=Ink]
多元极限 $\to$ 偏导与可微 $\to$ 链式与隐函数 $\to$ 梯度和方向导数 $\to$ Hessian与约束极值。
\end{tcolorbox}
\begin{examtip}考研主线：多元极限、偏导全微分、链式隐函数、方向导数和极值。常见题型：关系辨析、二阶求导、Hessian和乘子法。竞赛拓展：退化驻点、路径反例、紧区域上的全局优化。上述为学习任务分层，具体范围以所报考科目的当年大纲为准。\end{examtip}
\section{多元极限、偏导与全微分}\label{sec:6:1}
\subsection{概念、定理与推导}

\begin{definition}{多元极限、偏导与可微}
设$P=(a,b)$是定义域$D$的聚点。$\lim_{(x,y)\to P}f(x,y)=L$指：任意$\varepsilon>0$存在$\delta>0$，使$(x,y)\in D$且$0<\sqrt{(x-a)^2+(y-b)^2}<\delta$时有$|f(x,y)-L|<\varepsilon$；同一$\delta$控制全部方向。
在$P$为定义域内点时，偏导$f_x(P)=\lim_{h\to0}\frac{f(a+h,b)-f(a,b)}{h}$，$f_y$类似，极限均须为有限实数。
令$\rho=\sqrt{h^2+k^2}$。若存在常数$A,B$使
$f(a+h,b+k)-f(a,b)=Ah+Bk+o(\rho)$，称$f$在$P$可微；此时$A=f_x(P),B=f_y(P)$，全微分为$\dd f=f_x(P)\dd x+f_y(P)\dd y$。偏导只检查轴向，不能代替对全体增量的余项估计。
\end{definition}
\begin{theorem}{关系与充分条件}可微蕴含连续和偏导存在；偏导存在不保证连续。若偏导在一点邻域存在并在该点连续，则该点可微。\end{theorem}
\begin{proof}把总增量拆成先横向再纵向两段，分别用一元中值定理；偏导连续使系数误差为$o(1)$，乘以$h,k$后为$o(\rho)$。\end{proof}
极限否定只需两条路径；极限肯定则需统一估计，例如用$|xy|\le\frac{x^2+y^2}{2}$。混合偏导相等的常用充分条件是相邻二阶偏导连续。

\begin{method}{多元极限、偏导与全微分}条件：检查点处定义与邻域。步骤：否定用路径，证明用距离估计；可微直接查线性余项。易错点：所有直线路径相同仍不保证极限存在。\end{method}
\subsection{例题与方法比较}
\begin{example}{基础例题｜路径反例}{难度：$\star$}
\label{ex:6:1}\textbf{题目} 判断$\frac{xy}{x^2+y^2}$在原点极限。\par
\textbf{思路} 选轴线和对角线。\par
\textbf{解答} 方法一：沿$y=0$值0，沿$y=x$值$\frac{1}{2}$，故不存在。方法二：极坐标表达为$\cos\theta\sin\theta$，不随$r$衰减且依赖角度，不能有统一极限。\par
\textbf{方法点评} 路径法给最短否定，极坐标解释方向依赖。\par
\textbf{易错点或变式} 不能将$r\to0$与固定某个角度混同。
\end{example}
\begin{example}{考研风格例题｜统一估计}{难度：$\star\star$}
\label{ex:6:2}\textbf{题目} 把$f=\frac{x^2y}{x^2+y^2}$在原点补0，判断连续与可微。\par
\textbf{思路} 先估函数值，再估相对距离的余项。\par
\textbf{解答} $|f|\le|y|\to0$，故连续。轴上偏导均0。沿$y=x>0$，$f=\frac{x}{2}$，$\frac{f}{\rho}=\frac{1}{2\sqrt2}$，故不可微。另一验证：方向$\frac{1}{\sqrt2}(1,1)$的方向导数为$\frac{1}{2\sqrt2}$，与零梯度内积不符。\par
\textbf{方法点评} 连续与可微分别对应绝对误差和相对距离误差。\par
\textbf{易错点或变式} 偏导均为零不能推出可微。
\end{example}
\begin{example}{基础例题｜全微分}{难度：$\star$}
\label{ex:6:3}\textbf{题目} 求$f=x^2y+\sin y$在$(1,0)$的全微分。\par
\textbf{思路} 多项式与三角函数光滑。\par
\textbf{解答} 方法一：$f_x=2xy,f_y=x^2+\cos y$，故$\dd f=2\dd y$。方法二：增量$(1+h)^2k+\sin k=2k+2hk+h^2k+O(k^3)=2k+o(\rho)$。\par
\textbf{方法点评} 公式计算快，增量展开证明可微。\par
\textbf{易错点或变式} 线性主项中不含$h$不表示函数全局与$x$无关。
\end{example}
\section{复合函数与隐函数}\label{sec:6:2}
\subsection{概念、定理与推导}

\begin{theorem}{多元链式与隐函数，标准结论}可微复合映射的导数矩阵相乘。若$z=f(u,v)$且$u,v$依赖$x,y$，则$z_x=f_uu_x+f_vv_x$。若$F\in C^1$、$F(P)=0$且$F_z(P)\ne0$，在点附近唯一解出$z=z(x,y)$，$z_x=-\frac{F_x}{F_z},z_y=-\frac{F_y}{F_z}$。\end{theorem}
\begin{proof}将内层线性增量代入外层线性增量，余项仍为距离的小量。隐函数关系为\[F(x,y,z(x,y))=0.\]对该关系使用链式法则得到偏导公式；局部可解性来自隐函数定理。\end{proof}
方程组$F_1=F_2=0$解出$u,v$的充分条件为$\det\left(\frac{\partial(F_1,F_2)}{\partial(u,v)}\right)\ne0$；求导后解线性方程组。

\begin{method}{复合函数与隐函数}条件：区分自变量与中间变量。步骤：画依赖关系，每条路径导数相乘再相加；隐函数先整体求导。易错点：二阶求导时忘记系数也随变量变化。\end{method}
\subsection{例题与方法比较}
\begin{example}{考研风格例题｜链式}{难度：$\star\star$}
\label{ex:6:4}\textbf{题目} $z=u^2+v^2,u=x+y,v=x-y$，求$z_x$。\par
\textbf{思路} 链式或先展开。\par
\textbf{解答} 方法一：$z_x=2u+2v=4x$。方法二：$z=(x+y)^2+(x-y)^2=2x^2+2y^2$，直接得$4x$。\par
\textbf{方法点评} 先化简可消去中间变量，链式适合无显式表达的函数。\par
\textbf{易错点或变式} 偏导$x$时$y$固定。
\end{example}
\begin{example}{考研风格例题｜隐式二阶}{难度：$\star\star$}
\label{ex:6:5}\textbf{题目} 球面$x^2+y^2+z^2=1$上$z\ne0$，求$z_x,z_{xx}$。\par
\textbf{思路} 一阶关系再对$x$求导。\par
\textbf{解答} 方法一：$z_x=-\frac{x}{z}$，$z_{xx}=-\frac{1}{z}-\frac{x^2}{z^3}$。方法二：对$1+z_x^2+zz_{xx}=0$直接解出，得$z_{xx}=-\frac{1+z_x^2}{z}$，代回同式。\par
\textbf{方法点评} 隐式二次求导避免分式展开错误。\par
\textbf{易错点或变式} 不能把$z$当作常数求第二次导数。
\end{example}
\begin{example}{考研风格例题｜隐式方程组}{难度：$\star\star$}
\label{ex:6:6}\textbf{题目} $u+v=x,u-v=y$，求$u_x,v_y$。\par
\textbf{思路} 系数矩阵可逆。\par
\textbf{解答} 雅可比为$-2\ne0$。显式解$u=\frac{x+y}{2},v=\frac{x-y}{2}$，所以$u_x=\frac{1}{2},v_y=-\frac{1}{2}$。对两关系分别求偏导并解线性方程可得同样结果。\par
\textbf{方法点评} 线性例子显示一般隐函数组计算的结构。\par
\textbf{易错点或变式} 偏导变量变换时保持另一自变量不变。
\end{example}
\section{方向导数、梯度与切空间}\label{sec:6:3}

\begin{theorem}{空间曲线与曲面的切向}若曲线为$\bm r(t)$且$\bm r\prime(t_0)\ne0$，切线为$\bm r=\bm r(t_0)+s\bm r\prime(t_0)$，法平面为$\bm r\prime(t_0)\cdot[\bm r-\bm r(t_0)]=0$。若曲线由$F=G=0$给出且两梯度独立，其切向平行$\nabla F\times\nabla G$。\end{theorem}
\begin{proof}曲线的各分量满足微分关系。对$F(\bm r(t))=G(\bm r(t))=0$求导可知切向同时垂直于两梯度；独立性保证叉积非零。\end{proof}

\subsection{概念、定理与推导}

\begin{definition}{方向导数}沿单位向量$\bm l$的方向导数为$D_{\bm l}f(P)=\lim_{t\to0}\frac{f(P+t\bm l)-f(P)}{t}$。本册采用双侧定义。梯度为$\nabla f=(f_x,f_y,f_z)$。\end{definition}
\begin{theorem}{可微函数的方向导数}可微时$D_{\bm l}f=\nabla f\cdot\bm l$，若梯度非零，最大值$|\nabla f|$沿梯度方向达到，最小值$-|\nabla f|$沿反方向达到，梯度垂直于等值面；若梯度为零，所有单位方向的方向导数均为0。\end{theorem}
\begin{proof}沿$\bm l$代入线性增量；极值由Cauchy不等式，等值面的结论由沿曲线求导得到。\end{proof}
所有方向导数存在仍不保证可微；方向导数对方向线性是可微的必要特征之一。

\begin{method}{方向导数、梯度与切空间}条件：梯度公式要求可微且方向归一化。步骤：求梯度、单位化、内积；最快增长看梯度模。易错点：直接将非单位方向代入会多乘长度。\end{method}
\subsection{例题与方法比较}
\begin{example}{基础例题｜方向导数}{难度：$\star$}
\label{ex:6:7}\textbf{题目} $f=x^2+y^2$在$(1,2)$沿$(3,4)$方向的方向导数。\par
\textbf{思路} 先归一化。\par
\textbf{解答} $\nabla f=(2,4)$，单位方向$(\frac{3}{5},\frac{4}{5})$，方向导数$\frac{22}{5}$。直接代路径$(1+\frac{3t}{5},2+\frac{4t}{5})$，一阶$t$系数同样为$\frac{22}{5}$。\par
\textbf{方法点评} 定义计算可以检验归一化。\par
\textbf{易错点或变式} 直接内积得22是参数速度下的变化率。
\end{example}
\begin{example}{考研风格例题｜最快增长}{难度：$\star\star$}
\label{ex:6:8}\textbf{题目} $f=xyz$在$(1,1,2)$的最大方向导数。\par
\textbf{思路} 求梯度模。\par
\textbf{解答} 梯度$(2,2,1)$，模为3，最大方向为$\frac{1}{3}(2,2,1)$，最大方向导数3。\par
\textbf{方法点评} 线性优化转化为向量投影。\par
\textbf{易错点或变式} 梯度为零时所有方向一阶变化均为零，不能据此判极值。
\end{example}
\begin{example}{考研风格例题｜切平面}{难度：$\star\star$}
\label{ex:6:9}\textbf{题目} 求$z=x^2+y^2$在$(1,1,2)$处切平面。\par
\textbf{思路} 全微分或等值面梯度。\par
\textbf{解答} $z-2=2(x-1)+2(y-1)$，即$z=2x+2y-2$。令$F=x^2+y^2-z$，法向$(2,2,-1)$也给同式。\par
\textbf{方法点评} 图像形式与隐式形式一致。\par
\textbf{易错点或变式} 切平面只在接触点附近提供一阶近似。
\end{example}
\section{无条件极值与Taylor二次型}\label{sec:6:4}
\subsection{概念、定理与推导}

\begin{theorem}{多元Taylor与Hessian判别}若$f\in C^2$，则
\begin{equation}f(P+\bm h)=f(P)+\nabla f(P)\cdot\bm h+\dfrac12\bm h^TH(P)\bm h+o(|\bm h|^2).\end{equation}
驻点处Hessian正定给严格极小，负定给严格极大，不定给鞍点；半正定或半负定时不能仅凭二阶项决定。二元时$A=f_{xx},B=f_{xy},C=f_{yy}$，$AC-B^2>0$时按$A$符号判极值，$AC-B^2<0$时为鞍点。\end{theorem}
\begin{proof}沿线段设$g(t)=f(P+t\bm h)$，对$g$用一元Taylor。正定二次型在单位圆上有正的最小值，从而控制余项；不定二次型沿两方向异号。\end{proof}
紧区域上连续函数取得全局最值，候选点包括内部驻点、边界驻点、边界角点与不可微点。

\begin{method}{无条件极值与Taylor二次型}条件：内部驻点可用Hessian，边界须另算。步骤：解梯度零，分类二次型，比较所有候选值。易错点：判别式为0必须继续分析。\end{method}
\subsection{例题与方法比较}
\begin{example}{基础例题｜正定二次型}{难度：$\star$}
\label{ex:6:10}\textbf{题目} 求$f=x^2+xy+y^2-2x$的极小值。\par
\textbf{思路} 解线性驻点方程。\par
\textbf{解答} $2x+y=2,x+2y=0$给$(\frac{4}{3},-\frac{2}{3})$，$H=\begin{pmatrix}2&1\\1&2\end{pmatrix}$正定，最小值$-\frac{4}{3}$。平移后$f+\frac{4}{3}=u^2+uv+v^2=(u+\frac{v}{2})^2+\frac{3v^2}{4}\ge0$，也证明是全局最小。\par
\textbf{方法点评} Hessian判局部，配方直接给全局。\par
\textbf{易错点或变式} 驻点坐标代回需准确。
\end{example}
\begin{example}{考研风格例题｜鞍点}{难度：$\star\star$}
\label{ex:6:11}\textbf{题目} 判断$f=x^2-y^2$在原点的性质。\par
\textbf{思路} 两轴方向异号。\par
\textbf{解答} 沿$x$轴$f>0$，沿$y$轴$f<0$，故原点为鞍点。Hessian为$\operatorname{diag}(2,-2)$，行列式$-4<0$验证。\par
\textbf{方法点评} 路径比较比矩阵判别更直观。\par
\textbf{易错点或变式} 梯度零仅是驻点条件。
\end{example}
\begin{example}{考研风格例题｜退化判别}{难度：$\star\star$}
\label{ex:6:12}\textbf{题目} 判断$f=x^4+y^4$在原点的性质。\par
\textbf{思路} 二阶项全零，直接看高阶正性。\par
\textbf{解答} $f\ge0$，且仅原点取0，故严格全局极小。Hessian在原点为零矩阵，二阶判别失效却不影响直接结论。\par
\textbf{方法点评} 退化时回到定义或最低非零齐次项。\par
\textbf{易错点或变式} “判别法失效”不等于“无极值”。
\end{example}
\section{条件极值与竞赛优化}\label{sec:6:5}
\subsection{概念、定理与推导}

\begin{theorem}{Lagrange乘子，标准结论}设$f,g$在开集$U$上属于$C^1$，$P\in U$、$g(P)=0$，且$f$在约束集$\{g=0\}$上于$P$取得局部极值。若$\nabla g(P)\ne0$，则存在$\lambda$使$\nabla f(P)=\lambda\nabla g(P)$。对$m$个$C^1$等式约束$g_j=0$，要求$\nabla g_j(P)$线性无关，结论为$\nabla f(P)=\sum_{j=1}^m\lambda_j\nabla g_j(P)$。梯度退化、额外边界和不等式约束须另行检查。\end{theorem}
\begin{proof}可行曲线切向与$\nabla g$垂直，极值处$f$沿每个可行切向的变化率为零。因此$\nabla f$属于法空间，即约束梯度的张成空间。\end{proof}
\begin{extension}乘子给候选点而不自动判最值。紧约束可比较全部候选值；不紧约束要研究无穷远；约束梯度为零的奇异点需单独检查。不等式约束必须检查边界。\end{extension}

\begin{method}{条件极值与竞赛优化}条件：光滑正则约束。步骤：联立梯度比例与约束，补奇异点和边界，再比较函数值。易错点：把乘子条件当作充分条件。\end{method}
\subsection{例题与方法比较}
\begin{example}{竞赛风格例题｜圆上乘积}{难度：$\star\star\star$}
\label{ex:6:13}\textbf{题目} $x^2+y^2=1$，求$xy$最大最小值。\par
\textbf{思路} 乘子、平方恒等式或参数化。\par
\textbf{解答} 方法一：$(y,x)=2\lambda(x,y)$给$x=\pm y$，值为$\pm\frac{1}{2}$。方法二：$(x\mp y)^2\ge0$给$-\frac{1}{2}\le xy\le\frac{1}{2}$。方法三：$x=\cos t,y=\sin t$，$xy=\frac{\sin2t}{2}$。\par
\textbf{方法点评} 代数不等式最快，乘子法可推广，参数化降维。\par
\textbf{易错点或变式} 两种等号方向对应不同符号。
\end{example}
\begin{example}{竞赛风格例题｜固定和乘积}{难度：$\star\star\star$}
\label{ex:6:14}\textbf{题目} 正数$x+y+z=3$，证明$xyz\le1$。\par
\textbf{思路} 可用对数凹性或乘子。\par
\textbf{解答} 对数严格凹，$\frac13(\ln x+\ln y+\ln z)\le\ln(\frac{x+y+z}{3})=0$，得结论。另法：在闭单纯形考虑乘积，边界为0，内部乘子方程$yz=xz=xy$给$x=y=z=1$，最大值1。\par
\textbf{方法点评} 闭区域方法同时处理存在性与边界。\par
\textbf{易错点或变式} 使用对数必须先有正性。
\end{example}
\begin{example}{竞赛风格例题｜圆盘最值}{难度：$\star\star\star$}
\label{ex:6:15}\textbf{题目} 在$x^2+y^2\le1$上求$f=x+y$的最值。\par
\textbf{思路} 内部无驻点，边界是圆。\par
\textbf{解答} $\nabla f=(1,1)\ne0$，内部无极值。由Cauchy，$|x+y|\le\sqrt2$，在$\pm(\frac{1}{\sqrt2},\frac{1}{\sqrt2})$分别取最大最小。也可边界参数化为$\sqrt2\sin(t+\frac{\pi}{4})$。\par
\textbf{方法点评} 边界检查不可省略，简单线性目标用投影最有效。\par
\textbf{易错点或变式} 只解内部驻点会错误地得到“无最值”。
\end{example}
\begin{warning}偏导存在、所有方向导数存在、连续、可微是不同概念；使用梯度方向公式前确认可微；条件极值还须检查约束是否正则以及可行域边界。\end{warning}


\section{多元微分的严格条件与高阶方法}\label{mat:multivariable}
\begin{examtip}补充Euclid空间、紧致性、混合偏导、Jacobian、约束二次型与齐次函数。\end{examtip}
\subsection{邻域、开闭与紧致性}\label{mat:multivariable:topology}
\begin{method}{邻域、开闭与紧致性}\textbf{适用条件} 有限维Euclid空间。\par\textbf{实施与依据} 闭且有界等价于紧；连续像保紧，紧集上连续函数取最值。\end{method}
\begin{example}{专题例题｜邻域、开闭与紧致性}{难度：$\star\star$}
\textbf{题目} 解释$f(x,y)=\frac{1}{x^2+y^2}$在闭圆环$1\le x^2+y^2\le4$上最值。\par\textbf{思路与解答} 区域闭有界，函数连续；径向单调给最小$\frac{1}{4}$、最大1。若区域改为$0<x^2+y^2\le1$，失去闭性且函数无上界。\par\textbf{方法点评与边界} “闭集”必须相对于明确的空间定义。
\end{example}
\subsection{混合偏导交换的条件}\label{mat:multivariable:mixed}
\begin{method}{混合偏导交换的条件}\textbf{适用条件} 二阶混合偏导在邻域连续。\par\textbf{实施与依据} 对小矩形的四点差分别两次使用中值定理，再令边长趋零。\end{method}
\begin{example}{专题例题｜混合偏导交换的条件}{难度：$\star\star$}
\textbf{题目} 给出偏导存在但混合偏导不等的典型结构。\par\textbf{思路与解答} 令$f=\frac{xy(x^2-y^2)}{x^2+y^2}$，原点补0。沿坐标计算得$f_x(0,y)=-y$、$f_y(x,0)=x$，故$f_{xy}(0,0)=-1,f_{yx}(0,0)=1$。\par\textbf{方法点评与边界} 没有连续性假设时，形式上的交换求导可能错误。
\end{example}
\subsection{反函数与Jacobian链式}\label{mat:multivariable:jacobian}
\begin{method}{反函数与Jacobian链式}\textbf{适用条件} 映射$F\in C^1$且导数矩阵可逆。\par\textbf{实施与依据} 局部反函数满足$D(F^{-1})=(DF)^{-1}$；复合Jacobian相乘。\end{method}
\begin{example}{专题例题｜反函数与Jacobian链式}{难度：$\star\star$}
\textbf{题目} $u=x+y,v=x-y$，求反变换的Jacobian。\par\textbf{思路与解答} 正向行列式为$-2$，反向为$-\frac{1}{2}$。由$x=\frac{u+v}{2},y=\frac{u-v}{2}$直接计算也得到$-\frac{1}{2}$。\par\textbf{方法点评与边界} 面积换元用绝对值，导数矩阵本身不能擅自取绝对值。
\end{example}
\subsection{Euler齐次函数定理}\label{mat:multivariable:homogeneous}
\begin{method}{Euler齐次函数定理}\textbf{适用条件} $f(t\bm x)=t^mf(\bm x)$且可微。\par\textbf{实施与依据} 对$t$在1处求导，得到$\bm x\cdot\nabla f=mf$。\end{method}
\begin{example}{专题例题｜Euler齐次函数定理}{难度：$\star\star$}
\textbf{题目} 验证$f=(x^2+y^2)^{\frac{3}{2}}$的齐次恒等式。\par\textbf{思路与解答} 除原点外$f_x=3x\sqrt{x^2+y^2},f_y=3y\sqrt{x^2+y^2}$，故$xf_x+yf_y=3f$。原点两边均0，经定义也成立。\par\textbf{方法点评与边界} 只有可微及缩放关系同时成立时才可用。
\end{example}
\subsection{条件极值的二阶检验}\label{mat:multivariable:constrained2}
\begin{method}{条件极值的二阶检验}\textbf{适用条件} 约束正则，存在Lagrange乘子。\par\textbf{实施与依据} 将Lagrangian的Hessian限制在约束切空间，正定给严格条件极小。\end{method}
\begin{example}{专题例题｜条件极值的二阶检验}{难度：$\star\star$}
\textbf{题目} $x^2+y^2=1$上分析$f=x^2$。\par\textbf{思路与解答} 候选有$(0,\pm1)$与$(\pm1,0)$。前两点值0为条件最小，后两点值1为条件最大；参数$x=\cos t$给$f=\cos^2t$直接核验。\par\textbf{方法点评与边界} 限制在可行切向上，不要求整个空间Hessian正定。
\end{example}




\section{多元微分的反例、隐函数与极值深化}\label{deep:multivariable_source_deep}

\begin{examtip}从偏导、方向导数到可微逐级辨析，并将约束几何与Hessian判定联系。\end{examtip}

\subsection{极限、偏导与全微分的层级}\label{deep:multivariable_source_deep:differentiable}

\begin{method}{极限、偏导与全微分的层级}\textbf{适用条件} 多元极限须对所有趋近路径成立。\par\textbf{步骤与依据} 否定极限先选路径；证明可微则直接估计线性化余项除以距离。\par\textbf{易错点} 所有方向导数存在仍不足以推出可微。\end{method}

\begin{example}{竞赛风格｜所有方向导数存在仍不可微}{难度：$\star\star\star$}\label{deep:ex:multivariable_source_deep:directioncounter}\textbf{题目} 令$f(x,y)=\frac{x^3}{x^2+y^2}$（非原点），$f(0,0)=0$。求方向导数并判断可微。\par\textbf{思路} 沿直线算导数，再检查方向导数是否来自同一线性映射。\par\textbf{详细解答} 沿单位方向$(a,b)$，$\frac{f(ta,tb)}{t}=\frac{a^3}{a^2+b^2}=a^3$，故各方向导数存在。又$|f|\le|x|\le\sqrt{x^2+y^2}$，所以连续。偏导为$f_x(0)=1,f_y(0)=0$，若可微，方向导数应为$a$；而在$a=b=\frac{1}{\sqrt2}$时实际为$\frac{1}{2\sqrt2}$，故不可微。等价地沿$y=x$，$\frac{f(x,x)-x}{\sqrt{2x^2}}$不趋0。\par\textbf{方法点评} 本例准确展示“函数连续、所有方向导数存在，但不可微”，比只给不连续反例更强。\par\textbf{变式与易错点} 若方向导数不能表示成方向向量的同一个线性函数，就不可能全微分；仅逐方向计算存在性不够。\end{example}

\begin{example}{竞赛风格｜沿直线都趋零仍无极限}{难度：$\star\star\star$}\label{deep:ex:multivariable_source_deep:discontinuous}\textbf{题目} 令$g(x,y)=\frac{x^2y}{x^4+y^2}$，原点置0。检验直线和曲线趋近。\par\textbf{思路} 比较$y=kx$与$y=x^2$。\par\textbf{详细解答} 固定$k\ne0$，$g(x,kx)=\frac{kx}{x^2+k^2}\to0$；$k=0$时恒为0；竖直线$x=0$也为0。但沿$y=x^2$，$g=\frac{1}{2}$，所以原点不连续。说明所有直线方向的函数值极限一致仍不足以证明二元极限存在。\par\textbf{方法点评} 用非线性路径命中分母两项同阶的区域。\par\textbf{变式与易错点} 这里讨论沿线的函数极限，不是把所有方向导数也算作0。\end{example}

\subsection{隐函数、Jacobian与多元Taylor}\label{deep:multivariable_source_deep:implicitmany}

\begin{method}{隐函数、Jacobian与多元Taylor}\textbf{适用条件} 隐函数定理要求待解变量的Jacobian可逆。\par\textbf{步骤与依据} 把一阶偏导写成矩阵方程，二阶导数再求导；Taylor展开用Hessian二次型。\par\textbf{易错点} 偏导中保持不变的变量要明确，不能混用不同坐标系。\end{method}

\begin{example}{考研风格｜二元隐函数的一、二阶信息}{难度：$\star\star$}\label{deep:ex:multivariable_source_deep:implicitmatrix}\textbf{题目} $z=z(x,y)$由$x^2+y^2+z^2=1$在$(0,0,1)$附近定义，求二阶Taylor式。\par\textbf{思路} 先用$F_z=2z\ne0$确认上半球支。\par\textbf{详细解答} $z_x=-\frac{x}{z},z_y=-\frac{y}{z}$，在原点均0。由二次微分或$z=\sqrt{1-x^2-y^2}$得到$z_{xx}=z_{yy}=-1,z_{xy}=0$。因此$z=1-\frac12(x^2+y^2)+o(x^2+y^2)$。显式二项展开可进一步给余项$O((x^2+y^2)^2)$。\par\textbf{方法点评} 隐式法适用于不能显解的曲面，显式展开核对系数。\par\textbf{变式与易错点} 下半球支二阶项符号相反。\end{example}

\subsection{无条件、条件与退化临界点}\label{deep:multivariable_source_deep:extrema}

\begin{method}{无条件、条件与退化临界点}\textbf{适用条件} 内部极值用梯度零，边界用参数或乘数，最后比较候选值。\par\textbf{步骤与依据} 二阶正定给严格极小，半正定需高阶分析；约束下只看切空间上的二次型。\par\textbf{易错点} 拉格朗日方程给候选点，不自动给最大最小。\end{method}

\begin{example}{考研风格｜乘数、代入与不等式三解}{难度：$\star\star$}\label{deep:ex:multivariable_source_deep:lagrange}\textbf{题目} 在$x,y\ge0,x+y=2$下求$xy$的最大最小值。\par\textbf{思路} 三种方法都需检查端点。\par\textbf{详细解答} \textbf{乘数：}内部$(y,x)=\lambda(1,1)$，得$x=y=1$，值1；边界$(0,2),(2,0)$值0。\par\textbf{代入：}$xy=x(2-x)=1-(x-1)^2$，$0\le x\le2$，最大1、最小0。\par\textbf{不等式：}$xy\le[\frac{x+y}{2}]^2=1$，且非负，等号条件给全部极值点。\par\textbf{方法点评} 简单线性约束优先代入或不等式，乘数法适合更高维结构。\par\textbf{变式与易错点} 忘记非负约束时，最小值不存在。\end{example}

\begin{example}{竞赛风格｜Hessian退化的两个结局}{难度：$\star\star\star$}\label{deep:ex:multivariable_source_deep:degenerate}\textbf{题目} 比较$f=x^4+y^4$与$g=x^4-y^4$在原点的极值。\par\textbf{思路} 二阶导数全0，直接看高阶项符号。\par\textbf{详细解答} 两者原点梯度为0、Hessian都是零矩阵，二阶判别不定。$f\ge0$且只有原点为0，故严格极小；$g(x,0)>0,g(0,y)<0$，故鞍点。另$h=x^4+y^2$的Hessian半正定但仍严格极小，表明半正定不是否定结论。\par\textbf{方法点评} 退化时必须继续分析实际函数。\par\textbf{变式与易错点} 行列式为0只表示二阶判别无法决定。\end{example}

\section{图解：鞍点不是极值}
\begin{figure}[htbp]\centering
\begin{tikzpicture}\begin{axis}[width=.8\textwidth,height=6.5cm,view={45}{25},xlabel=$x$,ylabel=$y$,zlabel=$z$,domain=-1.2:1.2,y domain=-1.2:1.2,samples=17]
\addplot3[surf,shader=flat,draw=Ink!35,fill=Teal!35]{x^2-y^2};
\end{axis}\end{tikzpicture}
\caption{$f=x^2-y^2$在原点梯度为0，但沿$x$轴取正值、沿$y$轴取负值。Hessian为$\operatorname{diag}(2,-2)$，二次型不定，故原点是鞍点。可视化用于显示符号方向，严格判定依赖这两条截线或二次型。}\label{fig:review:saddle}
\end{figure}

\section{偏导条件的减弱、曲线信息与变量选择}\label{cmv:section}
\begin{examtip}
本节围绕三个判别问题展开：可微的充分条件能减弱到什么程度，沿曲线的信息能恢复多少偏导，以及同一个曲面上“对同一字母求导”为何可能得到不同结果。最后用线性微分算子和弯曲路径处理二阶方程与退化极值。
\end{examtip}

\subsection{一个偏导连续的可微判据}\label{cmv:onepartial}
\begin{theorem}{一个邻域偏导连续，另一个点处偏导存在}
设$f_x$在$P=(a,b)$的一个开矩形邻域内存在，并且$f_x$在$P$连续；若$f_y(P)$存在，则$f$在$P$可微。不要求$f_y$在整个邻域连续。
\end{theorem}
\begin{proof}
记$A=f_x(P),B=f_y(P)$，把增量沿两段坐标方向拆开：
\[
\begin{aligned}
f(a+h,b+k)-f(a,b)
={}&f(a+h,b+k)-f(a,b+k)\\
&+f(a,b+k)-f(a,b).
\end{aligned}
\]
第一段由关于$x$的一元中值定理等于
$h f_x(a+\vartheta h,b+k)$，其中$0<\vartheta<1$；
当$h=0$时该段为零。由$f_x$在$P$连续，第一段为$Ah+h\,o(1)$。
第二段由$f_y(P)$的定义等于$Bk+o(|k|)$。于是
\[
\Delta f=Ah+Bk+o\!\left(\sqrt{h^2+k^2}\right),
\]
因为$|h|,|k|\le\sqrt{h^2+k^2}$，故结论成立。
\end{proof}
\begin{method}{按连续的偏导选择拆增量的顺序}
\textbf{适用条件} 一个偏导在邻域存在且在考察点连续，另一个偏导只需在点处存在。\par
\textbf{步骤} 把需要使用中值定理的那一段安排给连续偏导；另一段沿过考察点的坐标轴，用导数定义。\par
\textbf{疑难点} “点处偏导存在”与“邻域内偏导存在”不同，不能把两者同时减弱为只在点处存在。
\end{method}
\begin{example}{考研风格｜偏导不连续仍可使用减弱判据}{难度：$\star\star$}
\label{cmv:ex:onepartial}
\textbf{题目} 令
\[
f(x,y)=
\begin{cases}
x^2+y^2\sin(\dfrac{1}{y}),&y\ne0,\\
x^2,&y=0.
\end{cases}
\]
判断$f$在原点的可微性，以及两个偏导的连续性。\par
\textbf{思路} 先检查连续的$x$偏导，再用原点处$y$偏导；也可直接估计余项。\par
\textbf{详细解答} 对所有点都有$f_x=2x$，所以$f_x$连续。原点处
$f_y(0,0)=\lim_{k\to0}k\sin(\frac{1}{k})=0$，由上述判据可微，$\dd f(0,0)=0$。
直接法给$|f(h,k)|\le h^2+k^2$，故$\frac{|f(h,k)|}{\sqrt{h^2+k^2}}\to0$，核对了结论。
当$y\ne0$时
\[
f_y(x,y)=2y\sin(\dfrac{1}{y})-\cos(\dfrac{1}{y}).
\]
取$y_n=(2\pi n)^{-1}$，有$f_y(0,y_n)=-1$，而$f_y(0,0)=0$，因此$f_y$在原点不连续。\par
\textbf{方法点评} 减弱判据与余项估计互相验证，并明确可微不要求全部偏导连续。\par
\textbf{易错点} 不能因一个偏导振荡就否定可微，也不能把充分条件误写为必要条件。
\end{example}

\subsection{两条曲线迹确定全微分的秩条件}\label{cmv:traces}
\begin{method}{把曲线切向导数写成线性方程组}
\textbf{适用条件} $f$在$P$可微，两条可微曲线$\gamma_1,\gamma_2$经过$P$，且切向量线性无关。\par
\textbf{步骤} 对已知的$f(\gamma_i(t))$求导，建立
$(f\circ\gamma_i)'=\nabla f(P)\cdot\gamma_i'$；
切向矩阵满秩时解出梯度。\par
\textbf{疑难点} 这只恢复点处的一阶信息，不能由两条曲线上的取值断定整个邻域的函数解析式。
\end{method}
\begin{example}{考研风格｜两条迹恢复微分，但不唯一恢复函数}{难度：$\star\star\star$}
\label{cmv:ex:traces}
\textbf{题目} $f$在$(1,1)$附近可微，且对充分小的$t$满足
\[
\begin{aligned}
f(1+t,\e^t)&=t(1+t)^2,\\
f(1+t,(1+t)^2)&=2(1+t)^2\ln(1+t).
\end{aligned}
\]
求$\dd f(1,1)$，并判断这些条件能否唯一确定$f$。\par
\textbf{思路} 两条曲线在该点的切向量分别为$(1,1)$、$(1,2)$。\par
\textbf{详细解答} 在$t=0$求导，记$A=f_x(1,1),B=f_y(1,1)$，得到
\[
A+B=1,\qquad A+2B=2.
\]
系数行列式为$1$，所以$A=0,B=1$，即$\dd f(1,1)=\dd y$。
函数$f_0(u,v)=u^2\ln v$满足两条迹关系。但对任意常数$c$，
\[
f_c(u,v)=u^2\ln v+
c\,[v-\e^{u-1}]\,[v-u^2]
\]
在$v>0$的小邻域中也满足相同关系：第一条曲线使第一个方括号为零，第二条曲线使第二个为零。因此函数并不唯一。\par
\textbf{方法点评} 曲线切向信息给一个有限维线性系统；两条迹不足以控制曲线之间的所有函数值。\par
\textbf{易错点} 若两条切向量平行，方程组可能只能确定一个方向导数。函数不唯一与全微分唯一并不矛盾。
\end{example}

\subsection{固定变量不同，偏导算子不同}\label{cmv:coordinates}
\begin{method}{先写清坐标图，再使用链式法则}
\textbf{适用条件} 同一约束曲面可以局部解出不同的因变量。\par
\textbf{步骤} 用下标说明保持不变的自变量；将目标函数拉回相应坐标平面后求导，或沿满足约束的切向量求导。\par
\textbf{疑难点} 曲面上$\frac{\partial}{\partial y}$必须与“固定$x$还是固定$z$”一起理解。
\end{method}
\begin{example}{考研风格｜同一曲面上两个关于$y$的偏导}{难度：$\star\star$}
\label{cmv:ex:coordinatepartials}
\textbf{题目} 在曲面
\[
x+y+z-4+\e^{-4}=\e^{-(x+y+z)}
\]
上考察$H=x^2yz$。在$P=(1,1,2)$，分别求保持$z$不变和保持$x$不变时$H$关于$y$的偏导。\par
\textbf{思路} 先识别约束只依赖$x+y+z$，再比较两个允许切向。\par
\textbf{详细解答} 令$s=x+y+z$，函数
$g(s)=s-4+\e^{-4}-\e^{-s}$满足$g(4)=0$且$g'(s)=1+\e^{-s}>0$。
所以该曲面恰为平面$x+y+z=4$。保持$z$不变时$x_y=-1$，故
\[
\left(\dfrac{\partial H}{\partial y}\right)_z
=H_y-H_x.
\]
保持$x$不变时$z_y=-1$，故
\[
\left(\dfrac{\partial H}{\partial y}\right)_x
=H_y-H_z.
\]
在$P$，$(H_x,H_y,H_z)=(4,2,1)$，两结果分别为$-2$和$1$。
直接代入$x=4-y-z$或$z=4-x-y$再求导，也得到相同结果。\par
\textbf{方法点评} 两个算子对应切向量$(-1,1,0)$和$(0,1,-1)$，几何方向不同，数值自然可以不同。\par
\textbf{易错点} 把空间函数$H_y$直接写成受约束的偏导，会漏掉随$y$变化的隐变量。
\end{example}

\subsection{二阶算子的特征坐标与弯曲路径}\label{cmv:operators}
\begin{method}{先分解微分算子，再展开二阶链式}
\textbf{适用条件} 常系数二阶偏微分方程，函数二阶连续可微。\par
\textbf{步骤} 把二次算子分解为一阶方向算子的乘积；寻找各方向的不变量，验证坐标变换可逆，再化成混合偏导方程。\par
\textbf{疑难点} 算子分解依赖混合偏导可交换；变换的Jacobian非零才保留二维坐标信息。
\end{method}
\begin{example}{竞赛风格｜把二阶方程化为混合偏导为零}{难度：$\star\star\star$}
\label{cmv:ex:operator}
\textbf{题目} 设$z\in C^2$满足$6z_{xx}+z_{xy}-z_{yy}=0$。取$u=x-2y,v=x+ay$，选择$a$使方程等价于$z_{uv}=0$，并写出局部通解。\par
\textbf{思路} 先观察$6D_x^2+D_xD_y-D_y^2=(3D_x-D_y)(2D_x+D_y)$。\par
\textbf{详细解答} 由链式法则，
\[
D_x=D_u+D_v,\qquad D_y=-2D_u+aD_v.
\]
于是$2D_x+D_y=(a+2)D_v$，
$3D_x-D_y=5D_u+(3-a)D_v$。取$a=3$，Jacobian
$\frac{\partial(u,v)}{\partial(x,y)}=5$非零，原方程化为$25z_{uv}=0$。
在一个$uv$开矩形上，$z_{uv}=0$使$z_u$与$v$无关，设$z_u=A'(u)$；
对$u$积分，得
\[
z(x,y)=A(x-2y)+B(x+3y),
\]
其中$A,B$为任意二阶连续可微的一元函数。代入两方向算子验证每一项分别被一个因子消去。\par
\textbf{方法点评} 分解法避免逐个展开全部二阶链式项；直接展开则是有效的系数核对。\par
\textbf{易错点} $a=-2$虽然会使某个系数消失，但此时$u=v$、Jacobian为零，不能当作有效换元。
\end{example}

\begin{example}{竞赛风格｜每条直线上极小仍不是二元极小}{难度：$\star\star\star$}
\label{cmv:ex:paraboliccorridor}
\textbf{题目} 令$f(x,y)=(y-x^2)(y-2x^2)$。证明原点是$f$沿每条过原点直线的严格局部极小点，却不是二元局部极小点。\par
\textbf{思路} 固定斜率后看最低阶项；否定二元极小则走进两条抛物线之间的窄带。\par
\textbf{详细解答} 沿$y=kx$，
\[
f(x,kx)=x^2(k-x)(k-2x).
\]
固定$k\ne0$时，取$0<|x|<\frac{|k|}{3}$，两括号同号，故$f>0$；$k=0$时$f=2x^4>0$。
竖直线$x=0$上$f=y^2>0$。因此每条固定直线都给严格局部极小。
但沿$y=\frac32x^2$有$f=-\frac{x^4}{4}<0$，而沿$y=0$有$f=2x^4>0$。
任意小邻域中都存在正、负函数值，原点为鞍点。\par
\textbf{方法点评} 固定方向时允许的邻域大小依赖斜率，不能把逐方向的结论升级为对所有方向统一的结论。\par
\textbf{疑难点} “每条直线”仍只检验了一类路径；此例的负值区宽度为$x^2$量级，直线检查会漏掉它。
\end{example}
\begin{figure}[htbp]\centering
\begin{tikzpicture}[x=2.7cm,y=2.0cm,>=stealth]
\fill[Rust!13] plot[domain=-.9:.9,samples=70] (\x,{\x*\x})
 -- plot[domain=.9:-.9,samples=70] (\x,{2*\x*\x}) -- cycle;
\draw[->,Ink](-1.03,0)--(1.08,0)node[right]{$x$};
\draw[->,Ink](0,-.08)--(0,1.84)node[above]{$y$};
\draw[Teal,thick]plot[domain=-.92:.92,samples=70](\x,{\x*\x});
\draw[Teal,thick]plot[domain=-.92:.92,samples=70](\x,{2*\x*\x});
\draw[Rust,dashed,thick]plot[domain=-.88:.88,samples=70](\x,{1.5*\x*\x});
\node[below left]at(0,0){$O$};
\node[right]at(.91,.83){$y=x^2$};
\node[right]at(.91,1.66){$y=2x^2$};
\node[Ink]at(0,1.43){窄带内$f<0$};
\node[Ink]at(-.65,.18){$f>0$};
\end{tikzpicture}
\caption{两条抛物线之间的负值带向原点收窄。沿中间抛物线趋近能检测到负值；每条固定直线上的局部正性无法替代二元极值定义。}
\label{cmv:fig:corridor}
\end{figure}

\section*{章末分层习题}\addcontentsline{toc}{section}{章末分层习题}
\refstepcounter{exercise}\label{q:6:1}\par\noindent\textbf{习题\theexercise\quad 基础题} 求$f=x^2y$的$f_x,f_y$。\quad\hyperref[a:6:1]{[答案]}\par
\refstepcounter{exercise}\label{q:6:2}\par\noindent\textbf{习题\theexercise\quad 基础题} 求$f=x^2+y^2$在$(0,0)$的全微分。\quad\hyperref[a:6:2]{[答案]}\par
\refstepcounter{exercise}\label{q:6:3}\par\noindent\textbf{习题\theexercise\quad 基础题} 求$xyz$的梯度。\quad\hyperref[a:6:3]{[答案]}\par
\refstepcounter{exercise}\label{q:6:4}\par\noindent\textbf{习题\theexercise\quad 考研风格} 求$z=\e^{x+y}$的$z_{xy}$。\quad\hyperref[a:6:4]{[答案]}\par
\refstepcounter{exercise}\label{q:6:5}\par\noindent\textbf{习题\theexercise\quad 考研风格} 判定$x^2+4y^2$在原点的极值。\quad\hyperref[a:6:5]{[答案]}\par
\refstepcounter{exercise}\label{q:6:6}\par\noindent\textbf{习题\theexercise\quad 考研风格} 求$f=x+y$沿$(1,1)$的单位方向导数。\quad\hyperref[a:6:6]{[答案]}\par
\refstepcounter{exercise}\label{q:6:7}\par\noindent\textbf{习题\theexercise\quad 竞赛拓展} $x^2+y^2=4$，求$x+2y$最大值。\quad\hyperref[a:6:7]{[答案]}\par
\refstepcounter{exercise}\label{q:6:8}\par\noindent\textbf{习题\theexercise\quad 竞赛拓展} 判断$x^2y^2$在原点是否严格极小。\quad\hyperref[a:6:8]{[答案]}\par
\refstepcounter{exercise}\label{q:6:9}\par\noindent\textbf{习题\theexercise\quad 竞赛拓展} 证明可微函数在内部局部极值处梯度为零。\quad\hyperref[a:6:9]{[答案]}\par
\section*{章末方法速查}\addcontentsline{toc}{section}{章末方法速查}
\par\noindent\begin{tabularx}{\textwidth}{p{23mm}XXX}\toprule 方法&适用条件&典型信号&易错点\\\midrule
距离估计&多元极限&$x^2+y^2$&只检查直线\\全微分&线性逼近&余项除距离&偏导代替可微\\Hessian&内部驻点&二次主项&退化误判\\乘子&正则约束&梯度平行&漏边界奇点\\
\bottomrule\end{tabularx}
